% !PRINTED 98
\weberpage{98}
%<<CONT

worin bei der Differentiation nicht Rücksicht genommen ist auf die Abhängigkeit $a_i^{(k)}=a_k^{(i)}$, dann wird die Formel (9)
\begin{equation}
 A\frac{\partial^2A}{\partial a_1^{(1)}\partial a_2^{(2)}}
 =\frac{\partial A}{\partial a_1^{(1)}}\frac{\partial A}{\partial a_2^{(2)}}
 -\left(\frac{\partial A}{\partial a_1^{(2)}}\right)^2.\tag{10}
\end{equation}

\begin{center}
{\large \S.\ 29.}

\vspace{0.5em}
{\large\gesperrt{Interpolation.}}
\end{center}

Wir haben schon in §. 9 unter einer besonderen Voraussetzung die Aufgabe gelöst, eine ganze rationale Function zu bestimmen, die für eine genügende Anzahl vorgeschriebener Werthe des Arguments gegebene Werthe annimmt. Wir wollen nun durch Anwendung der Determinanten diese Aufgabe allgemein lösen. Wir wollen aber folgende Bemerkungen vorausschicken.

Im §. 22, (12) ist der Werth der Determinante
\[
\begin{vmatrix}
 \alpha_1^{n-1}&\alpha_1^{n-2}&\cdots&\alpha_1&1\\
 \alpha_2^{n-1}&\alpha_2^{n-2}&\cdots&\alpha_2&1\\
 \cdots&\cdots&&\cdots&\cdots\\
 \alpha_n^{n-1}&\alpha_n^{n-2}&\cdots&\alpha_n&1
\end{vmatrix}
\]
gleich dem Differenzenproduct
\begin{equation}
(\alpha_1-\alpha_2)(\alpha_1-\alpha_3)\cdots(\alpha_1-\alpha_n)
(\alpha_2-\alpha_3)\cdots(\alpha_2-\alpha_n)\cdots(\alpha_{n-1}-\alpha_n)\tag{1}
\end{equation}
gefunden, und wir wollen jetzt zur Abkürzung dieses Differenzenproduct mit
\begin{equation}
 [\alpha_1,\alpha_2,\alpha_3,\ldots,\alpha_n]\tag{2}
\end{equation}
bezeichnen, so dass die Grösse $[\alpha_1,\alpha_2,\alpha_3,\ldots,\alpha_n]$ ihr Vorzeichen ändert, wenn zwei der Elemente $\alpha_1,\alpha_2,\ldots,\alpha_n$ vertauscht werden.

Wir definiren ferner eine ganze rationale Function $n$ten Grades $f(x)$ der Veränderlichen $x$ durch die Gleichung
\begin{equation}
 f(x)=(x-\alpha_1)(x-\alpha_2)\cdots(x-\alpha_n),\tag{3}
\end{equation}
so dass nach §. 13
\begin{align}
 f'(\alpha_1)&=(\alpha_1-\alpha_2)(\alpha_1-\alpha_3)\cdots(\alpha_1-\alpha_n),\notag\\
 f'(\alpha_2)&=(\alpha_2-\alpha_1)(\alpha_2-\alpha_3)\cdots(\alpha_2-\alpha_n),\tag{4}\\
 &\vdots\notag\\
 f'(\alpha_n)&=(\alpha_n-\alpha_1)(\alpha_n-\alpha_2)\cdots(\alpha_n-\alpha_{n-1}).\notag
\end{align}
%CONT>>
